\newpage
\section{Campagne à intensité turbulente variable et mise en évidence d'un
artefact de gravité en domaine périodique}
\label{sec:designA}

% ============================================================================
%  REECRIT le 2026-07-16. La campagne Design A a donne "pas de vitesse terminale"
%  a toutes les intensites ; le run LAMINAIRE de reference a revele que ce
%  comportement existe AUSSI en fluide au repos => ce n'etait pas la turbulence
%  mais un ARTEFACT NUMERIQUE (reduced.h incompatible avec la periodicite selon g).
%  Les vitesses d'ascension de TOUTE la campagne sont invalidees ; le correctif est
%  implemente et VALIDE (2026-07-22/24) sur quatre axes independants :
%    (1) convergence en maillage du cas laminaire a trois niveaux : u_inf = 10,93
%        (lvl6, delta/dx=0,87) / 12,274 (lvl7, 1,64) / 12,314 (lvl8, 3,28), soit
%        +0,32 % de lvl7 a lvl8 => le niveau 7 de production est converge ;
%    (2) courants parasites : u_max = 0,0093 a lvl7, soit 0,08 % de u_inf ;
%    (3) saut de Laplace vs 2 sigma/R = 255,7 : +1,8 % (lvl7), +0,48 % (lvl8) ;
%    (4) biais de sillage propre : -0,8 %, BORNE (le sillage sature), deux
%        methodes concordantes a 0,08 point.
%  Ne PAS reintroduire les conclusions "la turbulence supprime la remontee" :
%  c'etait l'epinglage sur la frontiere.
%
%  ATTENTION (correctif du 2026-07-22, event bubble de src/main.c) : jusqu'a cette
%  date le diagnostic ne sommait que sur les mailles taguees (f<0,5) et jetait la
%  moitie exterieure de la bande d'interface => volume sous-estime de ~4,6 points
%  et vitesse verticale biaisee de +0,35. Tout chiffre de volume "93-97 %" lu dans
%  un bubble.dat a 13 colonnes est un ARTEFACT DE MESURE : le VOF conserve a ~99 %.
% ============================================================================

Cette section rend rigoureuse la \og{}Série Intensité\fg{} esquissée plus haut :
on cherche à répondre à la question centrale du stage --- \emph{une turbulence
homogène isotrope accélère-t-elle ou ralentit-elle la remontée d'une bulle ?} ---
en faisant varier l'\textbf{intensité} de la turbulence, tous les autres
paramètres étant tenus fixes. Elle rapporte aussi, et surtout, la
\textbf{découverte d'un artefact numérique} qui invalide les vitesses
d'ascension mesurées jusqu'ici et dont la correction conditionne toute la suite.

\subsection{Cadre adimensionnel : trois nombres indépendants}

Le problème (à rapports de densité $\rho_1/\rho_2=850$ et de viscosité
$\mu_1/\mu_2=100$ fixés, géométrie fixée) fait intervenir six grandeurs
dimensionnelles : $g$, $\sigma$, $\nu$, $\varepsilon$ (dissipation turbulente),
$\rho$ et $D$. Le théorème de Vaschy--Buckingham ($6$ variables $-\,3$ dimensions
$M,L,T$) donne donc \textbf{exactement trois groupes sans dimension
indépendants}. On les choisit comme :
\begin{align}
    \text{Bond :}    \quad & Bo = \frac{(\rho_1-\rho_2)\,g\,D^2}{\sigma}, \\
    \text{Weber turbulent :} \quad & We_t = \frac{\rho\,\overline{\delta u^2}(D)\,D}{\sigma}
        \;\simeq\; \frac{2\rho\,\varepsilon^{2/3} D^{5/3}}{\sigma}, \\
    \text{Reynolds turbulent :} \quad & Re_\lambda = \frac{u'\lambda}{\nu}.
\end{align}
Le quatrième nombre usuel, le \textbf{Galilée} $Ga = \sqrt{g D^3}/\nu$ (ou de façon
équivalente le Morton $Mo = Bo^3/Ga^4$), n'est alors \emph{pas} libre : il est
\textbf{déterminé} par les trois autres. La définition de $We_t$ et le seuil de
rupture associé $We_{tC}=\mathcal{O}(1)$ suivent la revue de
\cite{legendre2025} (Éqs. 53--55).

À ces trois nombres correspondent trois boutons du solveur, réglés
indépendamment (arguments positionnels de \texttt{bubble}) :
\begin{center}
\renewcommand{\arraystretch}{1.3}
\begin{tabular}{lll}
    \toprule
    Bouton (arg.) & Règle & Effet de bord \\
    \midrule
    \texttt{BOND} & $Bo$ (exactement $Bo=$\,\texttt{BOND}) & --- \\
    \texttt{GRAV} $=|g|$ & $We_t$ à $Bo$ fixe ($We_t\propto 1/g$) & $Ga$ (dérivé) \\
    \texttt{NU} $=\nu$ & $Re_\lambda$ ($\nu\!\downarrow\Rightarrow Re_\lambda\!\uparrow$) & $Ga$ (dérivé) \\
    \texttt{KE\_TARGET} $\to\varepsilon$ & intensité réelle & $We_t$ \emph{et} $Re_\lambda$ \\
    \bottomrule
\end{tabular}
\end{center}

\paragraph{Choix du Design A.} On fait varier l'\textbf{intensité réelle} de la
turbulence via \texttt{KE\_TARGET}, à $Bo=1$ et $Ga\simeq 35$ \textbf{fixes}. Comme
un seul bouton ($\varepsilon$) pilote la turbulence, $We_t$ et $Re_\lambda$ montent
\textbf{ensemble} et ne peuvent être découplés : $We_t\propto\varepsilon^{2/3}$
(vite) et $Re_\lambda\propto\varepsilon^{1/6}$ (lentement).

\subsection{Protocole numérique : un précurseur par intensité}

Chaque intensité exige sa propre turbulence \emph{stationnaire}. Le pipeline est
donc à deux étages par point : (i) \textbf{précurseur} (\texttt{INJECT=0},
monophasique) qui reprend le champ \texttt{end\_we1} ($t=76$) et ré-équilibre la
turbulence à la nouvelle valeur de \texttt{KE\_TARGET} jusqu'à $t=160$ ;
(ii) \textbf{bulle} (\texttt{INJECT=1}) injectée dans cet état stationnaire et
laissée remonter jusqu'à $t=310$ ($150$ u.t.). Maillage adaptatif
(\texttt{maxlevel}~$=7$), critère \texttt{adapt\_wavelet} sur $\{f, |\omega|\}$,
plancher au niveau $5$.

\subsection{Résolution des échelles turbulentes}

Le tableau~\ref{tab:designA_resol} donne les échelles \textbf{mesurées} sur les
précurseurs équilibrés. Ces résultats restent \textbf{valides} : ils concernent la
phase monophasique, que l'artefact décrit plus bas n'affecte pas. Le critère
$k_{max}\,\eta > 1{,}5$ ($k_{max}=\pi\,2^7/L_0 = 3{,}35$) est vérifié partout, de
$7{,}7$ (we03) à $3{,}9$ (we20). Les $We_t$ réels collent aux cibles et
$Re_\lambda$ ne varie que de $18$ à $29$ (facteur ${\sim}1{,}6$) quand $We_t$ varie
d'un facteur~$6$ : le \textbf{verrouillage $We_t$--$Re_\lambda$} du Design A est
confirmé numériquement.

Point important pour le positionnement dans la littérature : nos précurseurs se
situent \textbf{dans le régime des DNS de référence}. \cite{legendre2025}
rapportent pour la DNS de Loisy \& Naso~\cite{loisy2017} $\eta/D\simeq0{,}098$, $\lambda/D\simeq1{,}0$
et $L_i/D\simeq2{,}1$ ; nous obtenons respectivement ${\simeq}0{,}09$--$0{,}11$,
${\simeq}0{,}9$--$1{,}1$ et ${\simeq}1{,}5$--$1{,}8$. La bulle est donc bien une
\textbf{grande bulle} ($\eta < D < L_i$), régime que la revue identifie comme une
\textbf{question ouverte}.

\begin{table}[htbp]
    \centering
    \renewcommand{\arraystretch}{1.3}
    \begin{tabular}{lcccccccccc}
        \toprule
        Run & $We_t$ & $k$ & $\varepsilon$ & $Re_\lambda$ & $\eta$ & $\lambda$ & $L$ & $\tau$ & $k_{max}\eta$ & $D/\eta$ \\
        \midrule
        we03 & 0.3 & 5 & 0.22 & 19 & 2.29 & 19.6 & 25 & 21 & 7.7 & 7.0 \\
        we07 & 0.7 & 11 & 0.59 & 25 & 1.79 & 17.5 & 29 & 17 & 6.0 & 8.9 \\
        we13 & 1.4 & 21 & 1.72 & 27 & 1.37 & 13.9 & 25 & 11 & 4.6 & 11.7 \\
        we16 & 1.5 & 25 & 2.03 & 29 & 1.31 & 13.9 & 27 & 11 & 4.4 & 12.2 \\
        we20 & 2.1 & 32 & 3.29 & 29 & 1.17 & 12.3 & 23 & 8 & 3.9 & 13.7 \\
        \bottomrule
    \end{tabular}
    \caption{Échelles turbulentes \emph{mesurées} des 5 précurseurs à l'équilibre
    ($\eta$ Kolmogorov, $\lambda$ Taylor, $L$ intégrale, $\tau=k/\varepsilon$ ;
    unités code). Résolution $k_{max}\eta>1{,}5$ partout. \textbf{Ces résultats
    ne sont pas affectés par l'artefact} (phase monophasique).}
    \label{tab:designA_resol}
\end{table}

\subsection{Un résultat suspect, puis sa réfutation par le cas laminaire}

\paragraph{Le constat initial.} Les cinq points ont tourné sans incident : aucun
NaN et \textbf{aucune fragmentation} --- une seule région gazeuse taguée sur toute
la fenêtre, là où une rupture ferait chuter brutalement le volume de la région mère
(cf.\ le run $KE=200$, $-88\%$). Le volume \emph{rapporté} à l'époque
($93{,}6$ à $96{,}5\%$ du théorique) sous-estimait toutefois systématiquement le
volume réel d'environ $4{,}6$ points, pour une raison purement instrumentale
détaillée plus bas (\S\ref{subsec:biais_volume}) : recalculé par sommes complètes,
il est conservé à $\ge97{,}1\%$ sur la campagne turbulente et à $\ge99\%$ sur le cas
laminaire. Mais la vitesse
de remontée nette, mesurée par régression de $z(t)$ au-delà du transitoire, est
\textbf{nulle à toutes les intensités} : $U_b\in[-0{,}22,\,+0{,}15]$, avec des
\emph{signes qui s'inversent} d'un point à l'autre. Le scénario est identique
partout (fig.~\ref{fig:designA_terminal}) : la bulle monte de ${\sim}50$ unités en
${\sim}20$ u.t., puis sa montée nette s'annule. L'interprétation tentante était
\og{}la turbulence supprime la remontée moyenne\fg{}.

\paragraph{Le test qui réfute.} Pour disposer de la vitesse de référence $u_\infty$
en fluide au repos --- indispensable au paramètre central du domaine
$\beta = u'/u_\infty$ \cite{legendre2025} --- un run \textbf{laminaire} a été
réalisé : même bulle, même $g$, mais liquide \emph{au repos} (vitesse annulée
partout à l'injection, forçage désactivé). Le champ reste effectivement quiescent
($k$ initial nul à la précision machine, $k\simeq10^{-3}$ en fin de run : le seul
sillage de la bulle). Or \textbf{la bulle s'y comporte exactement de la même
façon} (tab.~\ref{tab:lam_artefact}) : elle accélère jusqu'à $u_z\simeq12$, monte
$53{,}4$ unités, puis $u_z$ \emph{devient négatif} et elle se \textbf{fige} pour
les $140$ u.t. restantes.

\begin{table}[htbp]
    \centering
    \renewcommand{\arraystretch}{1.3}
    \begin{tabular}{lcccccc}
        \toprule
        $t$ & 162 & 165 & 170 & 175 & 190 & 310 \\
        \midrule
        montée $z(t)-z_0$ & 10.6 & 42.5 & 54.2 & 53.6 & 53.3 & \textbf{53.4} \\
        $u_z$ & $+9{,}0$ & $+11{,}9$ & $\mathbf{-3{,}6}$ & $-1{,}4$ & $-0{,}1$ & $\mathbf{+0{,}02}$ \\
        \bottomrule
    \end{tabular}
    \caption{Run \textbf{laminaire} (liquide au repos, $g=4$, $Bo=1$). Une bulle
    flottante en fluide quiescent doit monter indéfiniment à sa vitesse terminale.
    Elle s'arrête : le comportement n'est donc \emph{pas} dû à la turbulence.}
    \label{tab:lam_artefact}
\end{table}

\paragraph{Diagnostic : la gravité réduite est incompatible avec la périodicité
selon $\mathbf{g}$.} Le module \texttt{reduced.h} exprime la poussée non pas comme
une force volumique mais comme une force interfaciale dérivant du potentiel
\begin{equation}
    \phi = [\rho]\;\mathbf{G}\cdot(\mathbf{x}-\mathbf{Z}),
    \label{eq:phi_reduced}
\end{equation}
\textbf{linéaire en $z$}, donc \textbf{non périodique} : il saute de $[\rho]\,g\,L_0$
à la traversée de la frontière $z=L_0$. Dès que la bulle atteint le haut du
domaine, une partie de son interface est à $z\lesssim L_0$ (où $\phi$ est maximal)
et l'autre se replie en $z\gtrsim0$ (où $\phi$ est nul) : ce saut engendre une
force parasite considérable qui \textbf{épingle la bulle sur la frontière}.

La signature est sans appel : la bulle est injectée en $z_c=60$ et se fige en
$z_c=113{,}4$, alors que $L_0 - R = 120 - 8 = 112$. Elle s'arrête donc
\emph{exactement} lorsque son sommet atteint la frontière périodique. Et
\textbf{tous} les runs du projet --- $We\!\approx\!1$ (lvl7 et lvl8), les cinq
points Design A, le laminaire --- se figent à cette même position, quelle que soit
l'intensité turbulente. Ce n'est pas de la physique, c'est de la géométrie.

\begin{figure}[htbp]
    \centering
    \includegraphics[width=\textwidth]{designA_terminal.png}
    \caption{\textbf{Signature de l'artefact.} À gauche : moyenne cumulée
    $U_{cum}(t)$, qui décroît vers $0$ sans jamais plateauer. À droite :
    trajectoire verticale nette --- toutes les bulles montent ${\sim}50$ unités
    puis stagnent ou sont ballottées \emph{autour de la même altitude}
    $z_c\simeq L_0-R$. Le cas laminaire (non représenté ici) reproduit ce blocage
    en l'absence totale de turbulence, ce qui établit l'origine numérique du
    phénomène.}
    \label{fig:designA_terminal}
\end{figure}

\subsection{Portée de l'artefact : ce qui tombe, ce qui reste}

\begin{center}
\renewcommand{\arraystretch}{1.3}
\begin{tabular}{p{0.52\textwidth}l}
    \toprule
    Résultat & Statut \\
    \midrule
    Vitesses d'ascension $U_b$, \og{}absence de vitesse terminale\fg{},
    \og{}la turbulence supprime la remontée\fg{} & \textbf{invalidés} (artefact) \\
    Convergence en maillage lvl7$\approx$lvl8 sur $U_b$ & invalidée (les deux
    s'épinglaient au même endroit) --- \textbf{refaite depuis} sur le cas laminaire
    corrigé, cf.~\S\ref{subsec:convergence} \\
    Coefficient de traînée $C_d$ & sans objet ($U_b$ non physique) \\
    Précurseurs, caractérisation de la turbulence, tab.~\ref{tab:designA_resol} &
    \textbf{valides} ($g$ n'agit pas en monophasique) \\
    Non-fragmentation jusqu'à $We_t\simeq2{,}1$ & \textbf{à reconfirmer} (la bulle
    était maintenue dans un champ de force parasite) \\
    \bottomrule
\end{tabular}
\end{center}

\subsection{Correctif : une poussée compatible avec la périodicité}

On abandonne \texttt{reduced.h} au profit d'une \textbf{force volumique
périodique}, obtenue en retranchant la densité moyenne $\rho_m$ du domaine :
\begin{equation}
    \rho\,\mathbf{a} = (\rho - \rho_m)\,\mathbf{g}
    \qquad\Longleftrightarrow\qquad
    a_z = -g\left(1 - \frac{\rho_m}{\rho}\right).
    \label{eq:fix}
\end{equation}
Ses trois propriétés répondent exactement au problème :
\begin{enumerate}
    \item $\rho(\mathbf{x})$ est périodique, donc la force l'est aussi : \textbf{aucun
    potentiel linéaire}, la bulle traverse la frontière et \emph{continue} de monter ;
    \item sa moyenne est nulle, $\int(\rho-\rho_m)\,\mathbf{g}\,\mathrm{d}V = 0$ :
    pas d'emballement de la quantité de mouvement totale, ce qu'exige un domaine
    triplement périodique ;
    \item en monophasique ($f=1$) on a $\rho=\rho_m$ donc $a_z=0$ \emph{exactement} :
    le précurseur n'est pas affecté par $g$. \textbf{Les précurseurs déjà calculés
    restent donc valides et réutilisables} --- seuls les runs bulle sont à refaire.
\end{enumerate}
Contrepartie assumée : la formulation n'est plus \emph{well-balanced} (la bulle
reçoit $a_z\simeq g\,\rho_m/\rho_2$, équilibré par le gradient de pression via la
projection, l'accélération nette valant ${\sim}2g$ du fait de la masse ajoutée).
C'est le prix de la compatibilité périodique, et la formulation standard des DNS
de bulles en boîte périodique.

\paragraph{Validation : le test réfute l'artefact.} Le test décisif est le
\emph{même} cas laminaire relancé avec \eqref{eq:fix}. Le contraste est sans
ambiguïté (tab.~\ref{tab:fix_verdict}) : au lieu de se figer à $53{,}4$, la
trajectoire dépliée croît \textbf{linéairement} ($R^2 = 0{,}9999$) sur
$1889$ unités, soit \textbf{15,7 traversées} du domaine, et la vitesse se stabilise
à $u_\infty = 12{,}6$ ($u_z = 13{,}1 \pm 0{,}5$). La bulle franchit la frontière
périodique et poursuit son ascension : le comportement attendu d'une bulle
flottante en fluide au repos est restauré.

\begin{table}[htbp]
    \centering
    \renewcommand{\arraystretch}{1.3}
    \begin{tabular}{lcc}
        \toprule
        Cas laminaire ($g=4$, $Bo=1$, $150$ u.t.) & \texttt{reduced.h} & Correctif \eqref{eq:fix} \\
        \midrule
        Montée nette (unités) & $53{,}4$ puis \textbf{figée} & $\mathbf{1889}$ \\
        Traversées de domaine & $0$ & $\mathbf{15{,}7}$ \\
        $u_z$ en fin de run & $0{,}02$ & $\mathbf{13{,}6}$ \\
        Linéarité de $z(t)$ & --- & $R^2 = 0{,}9999$ \\
        $u_\infty$ & non définie & $\mathbf{12{,}6}$ \\
        \bottomrule
    \end{tabular}
    \caption{Verdict du correctif sur le cas limite laminaire.}
    \label{tab:fix_verdict}
\end{table}

\paragraph{Validation croisée avec la littérature.} La bulle corrigée est
ellipsoïdale et son taux de déformation est \emph{prédictible}. Cette première
mesure ($u_\infty\simeq12{,}6$) est toutefois entachée par la dérive de volume
signalée ci-dessous ; la confrontation quantitative à la théorie est donc reportée
au run définitif à repère mobile (§\ref{subsec:theorie}), où le diamètre reste
constant.

\subsection{Un défaut que l'artefact masquait : la conservation du volume}
\label{subsec:biais_volume}

\paragraph{Avertissement préalable : deux effets à ne pas confondre.} Les chiffres
de volume figurant dans les versions antérieures de ce document mélangeaient un
\emph{biais de mesure} et une \emph{dérive physique}. Il faut les séparer.

\begin{itemize}
  \item \textbf{Un biais de mesure, constant, d'environ $4{,}6$ points.} Jusqu'au
  \nobreak22 juillet 2026, l'\texttt{event bubble} du solveur construisait un
  marqueur binaire $m=(f<0{,}5)$, le segmentait en régions connexes, puis ne sommait
  que sur la région retenue. Or une maille à $f=0{,}6$ contient $40\%$ de gaz mais
  n'est pas marquée : son gaz était \emph{jeté}. La bande d'interface faisant une à
  deux mailles, cela revenait à écarter systématiquement sa moitié extérieure.
  Mesuré sur \texttt{snapshot-200.00} du run laminaire : $806$ mailles ignorées,
  $2038{,}1$ au lieu de $2134{,}7$, soit $4{,}5\%$ du volume. Le même défaut biaisait
  la vitesse verticale de $+0{,}35$ ($2{,}8\%$ de $u_\infty$), et dans ce sens précis
  parce qu'en repère mobile le liquide défile vers le bas : les mailles écartées,
  les plus riches en liquide, sont les plus négatives, et les retirer relève la
  moyenne. Tout volume \og{}$93$--$97\%$\fg{} lu dans un \texttt{bubble.dat}
  à $13$ colonnes relève de cet artefact.
  \item \textbf{Une dérive d'advection VOF, réelle et croissante.} Elle n'apparaît
  que lorsque la bulle traverse effectivement le maillage, c'est-à-dire dans les runs
  \emph{sans} repère mobile.
\end{itemize}

\paragraph{La dérive, une fois la mesure corrigée.} Le correctif de la poussée
révèle donc un second problème, jusque-là \textbf{invisible} parce que la bulle
épinglée ne parcourait que $53$ unités. Libérée, elle traverse le domaine seize fois
en $150$ u.t., et son volume \textbf{croît de $\mathbf{98\%}$ à $\mathbf{122\%}$}
du théorique (lecture brute $93{,}6\to116{,}1\%$, corrigée du biais ci-dessus).
Le diagnostic est net : la croissance est \emph{lisse} en fonction du temps
($R^2 = 0{,}99$), et non par paliers ; elle ne se produit donc pas aux traversées de
frontière mais s'accumule continûment, au rythme de $\mathbf{1{,}2\%}$
\textbf{par traversée de domaine} ($+12{,}5\%$ par $100$~u.t.). C'est une
\textbf{erreur d'advection VOF}.

Ce défaut préexistait : la bulle épinglée ne parcourait que $53$ unités, contre
$1889$ à présent --- $35$ fois plus d'advection à travers le maillage. L'artefact
le dissimulait. Il doit être résolu avant toute exploitation quantitative : un
volume croissant est un diamètre croissant, donc une vitesse terminale qui dérive
(on observe effectivement $u_z$ passer de $12{,}5$ à $13{,}7$).

\subsection{Version actuelle du code : le repère mobile}

La cause étant le \emph{déplacement de l'interface à travers la grille}, le remède
est de supprimer ce déplacement plutôt que de raffiner indéfiniment. On exploite
l'\textbf{invariance galiléenne} : retrancher une vitesse uniforme à tout le champ
ne modifie pas la physique. On suit donc la bulle --- elle reste au centre du
domaine, le liquide défile devant elle :
\begin{equation}
    u_z \leftarrow u_z - \delta, \qquad
    \delta = \frac{\mathrm{d}t}{\tau_f}\,\overline{u_z}\big|_{\text{gaz}},
    \qquad U_f \leftarrow U_f + \delta ,
    \label{eq:frame}
\end{equation}
$U_f$ étant la vitesse cumulée du repère. Aucune information n'est perdue : la
trajectoire de laboratoire est $z_{\text{lab}}(t) = \int U_f\,\mathrm{d}t$ et la
vitesse terminale est $U_f$ à l'équilibre (fichier \texttt{frame.dat}).

Trois points méritent d'être explicités :
\begin{itemize}
    \item \textbf{Pourquoi cela corrige le volume} : l'interface ne se déplace plus
    dans le maillage, donc l'erreur d'advection VOF disparaît par construction.
    \item \textbf{Pourquoi une relaxation} ($\tau_f = 1$~u.t.) plutôt qu'un recalage
    brutal : un repère qui \emph{accélère} est non inertiel et introduit une
    pseudo-force. À la vitesse terminale $U_f$ est constant, le repère est donc
    inertiel et \eqref{eq:frame} est \emph{exact} ; la relaxation évite une
    pseudo-force violente pendant le transitoire et préserve les fluctuations
    turbulentes de la bulle. À l'équilibre $\overline{u_z}|_{\text{gaz}}\to0$ donc
    $\delta\to0$ : pas de dérive parasite.
    \item \textbf{Pourquoi conserver \eqref{eq:fix} plutôt que revenir à
    \texttt{reduced.h}} : le repère mobile maintenant la bulle loin des frontières,
    \texttt{reduced.h} redeviendrait techniquement utilisable (et bénéficierait de
    son caractère well-balanced). Mais si le suivi décrochait --- en turbulence
    intense, par exemple --- la bulle atteindrait la frontière et l'épinglage
    reviendrait \emph{silencieusement}. La force volumique \eqref{eq:fix} est
    robuste quelle que soit la position de la bulle ; c'est ce filet de sécurité qui
    a été privilégié.
\end{itemize}

\paragraph{Verdict (run lvl7, $100$ u.t.).} Le test de cette version est, de
nouveau, le \emph{même} cas laminaire ($g=4$, $Bo=1$), sur $t\in[160,260]$. Les
trois attentes sont satisfaites simultanément (tab.~\ref{tab:frame_verdict}) :
\begin{itemize}
    \item la vitesse d'ascension se stabilise sur un \textbf{plateau net}
    $u_\infty = 12{,}27$ (pente résiduelle $-2\times10^{-3}$/u.t. sur $[210,260]$),
    lu directement dans \texttt{frame.dat} ; la bulle parcourt $z_{\text{lab}}=1219$
    unités, soit $10{,}2$ traversées-équivalent de domaine ;
    \item le \textbf{volume est conservé}. Mesuré correctement --- somme complète
    pondérée $\int(1-f)\,\mathrm{d}V$ sur tout le domaine, et non sur les seules
    mailles marquées (\S\ref{subsec:biais_volume}) --- il vaut $99{,}05\%$ du
    théorique ($V_{\text{th}}=2144{,}7$) à $t=161$ et $99{,}88\%$ à $t=260$, soit une
    dérive résiduelle de $\mathbf{+0{,}8\%}$ sur $100$~u.t. ($+0{,}6\%$ par
    régression), contre $+12{,}5\%$ sans repère mobile. Le diamètre équivalent passe
    de $15{,}95$ à $15{,}99$ pour $D=16$ : \textbf{constant à $0{,}3\%$ près}. Le
    \og{}déficit initial de $6\%$\fg{} rapporté antérieurement n'existe pas --- c'était
    le biais du diagnostic ; la morsure réelle de \texttt{fraction()} à l'injection
    vaut $0{,}9\%$ ;
    \item le \textbf{garde-fou de quantité de mouvement} (fichier
    \texttt{momentum.dat}, ajouté pour cette campagne) confirme l'absence de force
    parasite : en repère laboratoire $|p_x|,|p_y|\le4\times10^{-3}$ et
    $|p_{z,\text{lab}}|\le0{,}04$, soit $0{,}35\%$ de $u_\infty$ sur tout le run.
\end{itemize}

\begin{table}[htbp]
    \centering
    \renewcommand{\arraystretch}{1.3}
    \begin{tabular}{lcc}
        \toprule
        Cas laminaire ($g=4$, $Bo=1$) & Force volumique seule & \textbf{+ repère mobile} \\
        \midrule
        $u_\infty$ & $12{,}6$ (biaisé) & $\mathbf{12{,}27}$ (plateau) \\
        Dérive de volume /$100$ u.t. & $+12{,}5\%$ & $\mathbf{+0{,}8\%}$ \\
        Déplacement dans la grille & $1889$ u. ($16$ trav.) & ${\sim}0$ (bulle centrée) \\
        qdm lab conservée ($p_{z,\text{lab}}/u_\infty$) & non diagnostiquée & $\mathbf{0{,}35\%}$ \\
        \bottomrule
    \end{tabular}
    \caption{Apport du repère mobile sur le cas limite laminaire. La physique
    ($u_\infty$) est préservée ; le volume cesse de dériver, ce qui rend $u_\infty$
    fiable (diamètre constant).}
    \label{tab:frame_verdict}
\end{table}

\subsection{Confrontation à la théorie}
\label{subsec:theorie}

Avec la valeur \emph{fiable} $u_\infty = 12{,}27$ (diamètre constant), la bulle
laminaire se caractérise par
\begin{equation}
    Re_b = \frac{u_\infty D}{\nu} = 107, \quad
    Fr = \frac{u_\infty}{\sqrt{gD}} = 1{,}53, \quad
    We = \frac{\rho\,u_\infty^2 D}{\sigma} = 2{,}34, \quad
    Ga = 70, \quad Mo = 4{,}1\times10^{-8}.
\end{equation}

\paragraph{Coefficient de traînée.} À l'équilibre poussée--traînée, le coefficient
mesuré vaut
\begin{equation}
    C_D = \frac{4}{3}\,\frac{(\rho_1-\rho_2)\,g\,D}{\rho_1\,u_\infty^2} = 0{,}57.
\end{equation}
Il se situe \textbf{entre} les deux références classiques à $Re_b = 107$ : la
corrélation de Mei \emph{et al.}~\cite{mei1994} pour une bulle \emph{propre} (interface
mobile, $\lambda=\mu_2/\mu_1=0{,}01\ll1$), $C_D^{\text{propre}} = 0{,}35$, et celle
de Schiller--Naumann pour une sphère \emph{solide} (interface rigide),
$C_D^{\text{solide}} = 1{,}06$. L'excès de $62\%$ sur la bulle propre s'explique par
deux effets qui augmentent tous deux la traînée : la \textbf{déformation oblate}
($\chi = 1{,}38$, section frontale accrue) et le caractère \textbf{instationnaire du
sillage} (trajectoire oblique, ci-dessous). L'ordre de grandeur et le
positionnement entre les deux bornes sont donc physiquement cohérents.

\paragraph{Déformation.} La corrélation de \cite{legendre2025} (Éq.~45),
$\chi = [1 - (9/64)\,We\,(1+0{,}2\,Mo^{1/10})^{-1}]^{-1}$, prédit
$\chi_{\text{préd}} = 1{,}47$ pour nos nombres, à comparer au
$\chi_{\text{mes}} = 1{,}38$ mesuré : accord à $6\%$, \emph{sans aucun ajustement}.

\subsection{La trajectoire oblique : une instabilité de sillage, pas un artefact}

Une dérive horizontale subsiste, mais elle est de nature \textbf{différente} de la
\og{}diagonale parfaite\fg{} du run sans repère : celle-ci croissait
\emph{linéairement et sans borne} le long de l'axe \emph{exact} de la grille $(1,1)$
(azimut $45^\circ$) --- signature d'une \textbf{erreur d'advection} systématique, désormais
supprimée par le repère mobile. Ici, la dérive s'oriente selon un azimut
${\simeq}66^\circ$ (rapport $\mathrm{d}y/\mathrm{d}x\simeq2{,}2$), \emph{non aligné} sur
un axe du maillage, et sa vitesse reste modérée ($|u_h|$ de $5$ à $18\%$ de
$u_\infty$ le long du run, moyenne ${\simeq}11\%$). Surtout, le garde-fou
\texttt{momentum.dat} montre que la quantité de mouvement horizontale est
\textbf{conservée} ($|p_x|,|p_y|\le4\times10^{-3}$) : cette dérive n'est donc pas due
à une force parasite, c'est une \emph{redistribution interne} bulle/liquide.

L'interprétation la plus probable est une \textbf{instabilité de trajectoire} du
sillage : à $Bo=1$, le seuil de perte de la trajectoire rectiligne est
$Ga_c\simeq44$ (diagramme $(Ga,Bo)$ de Cano-Lozano \& Magnaudet~\cite{canolozano2016}), et avec
$Ga=70>Ga_c$ la bulle doit adopter un chemin oblique.

\paragraph{Réserve levée.} Une version antérieure de ce document laissait ouverte
la possibilité d'une \og{}composante numérique résiduelle\fg{} dans cette dérive.
Deux mesures la ferment.

\emph{(i) Les courants parasites sont négligeables.} Un cas dédié a été construit
(\texttt{src/static\_bubble.c}) : une bulle sphérique au repos dans un liquide au
repos, \textbf{sans gravité et sans forçage}, configuration dont la solution exacte
est $\mathbf{u\equiv0}$ --- la tension de surface y est exactement équilibrée par le
saut de Laplace. Toute vitesse mesurée est donc de l'erreur numérique pure. On
obtient $u_{\max}=9{,}3\times10^{-3}$ à \texttt{lvl7}, soit $\mathbf{0{,}08\%}$ de
$u_\infty$ : deux ordres de grandeur sous la dérive latérale observée. Le saut de
pression numérique vaut $260{,}3$ contre $2\sigma/R=255{,}7$ attendu, soit $+1{,}8\%$
($+0{,}48\%$ à \texttt{lvl8}), ce qui valide au passage l'implémentation de la
tension de surface.

\emph{(ii) La dérive latérale sature.} $|u_h|$ ne croît pas indéfiniment : elle
s'installe en \textbf{cycle limite} borné, d'amplitude $[0{,}14;2{,}57]$ et de
période ${\simeq}16{,}2$~u.t. --- exactement celle du lâcher tourbillonnaire
($St = fD/u_\infty = 0{,}080$, dans la fourchette $0{,}05$--$0{,}15$ de la
littérature). Une dérive numérique croîtrait sans borne ; un cycle limite à la
fréquence du sillage est une signature physique.

Point crucial, inchangé : cette dérive est \emph{perpendiculaire} à l'ascension et
\textbf{n'affecte pas} la mesure de $u_\infty$.

\subsection{Point de vigilance : l'interaction avec le sillage propre}

Un dernier effet mérite d'être signalé --- et, cette fois, d'être \emph{chiffré}
plutôt qu'écarté. En boîte triplement périodique, le sillage laissé par la bulle
\textbf{se reboucle} et la bulle finit par le re-rencontrer : sur $100$ u.t. et
${\sim}10$ traversées, l'énergie cinétique de fluctuation du liquide passe de
$k\simeq0{,}05$ à $0{,}45$, alors qu'elle devrait rester nulle en fluide au repos.
La \og{}référence en liquide quiescent\fg{} n'en est donc plus tout à fait une.

\paragraph{Le biais est borné, et il vaut $-0{,}8\%$.} Deux points le montrent.
D'abord, l'accumulation \textbf{sature} : $k$ atteint $0{,}45$ vers $t\simeq190$ puis
plafonne, le sillage rebouclé trouvant un équilibre production/dissipation. Le biais
n'est donc pas cumulatif, et allonger le run ne l'aggraverait pas. Ensuite, son
amplitude se mesure de \textbf{deux façons indépendantes qui concordent} :
\begin{itemize}
    \item en traitant le sillage propre comme une turbulence de fond d'intensité
    $\beta_{\text{sillage}}=u'_{\text{sillage}}/u_\infty=0{,}045$ et en l'injectant
    dans la loi petit-$\beta$ ajustée sur nos propres points
    (\S\ref{subsec:resultats_turb}), on prédit $-0{,}76\%$ ;
    \item la dérive de $u_\infty$ directement mesurée sur $[220,260]$ vaut
    $-0{,}84\%$.
\end{itemize}
Les deux méthodes s'accordent à $0{,}08$ point. Le biais de sillage est donc
\textbf{un à deux ordres de grandeur sous le signal étudié} ($-6$ à $-39\%$) :
c'est une limite à assumer explicitement, non un obstacle. Sur les points turbulents,
l'ajouter en quadrature ($\beta_{\text{eff}}=\sqrt{\beta^2+\beta_{\text{sillage}}^2}$)
déplace $\beta$ de $+4{,}4\%$ au point le plus faible à $+0{,}7\%$ au plus fort.

Trois parades restent envisageables si l'on voulait le supprimer plutôt que le
borner, par ordre de préférence :
\begin{enumerate}
    \item \textbf{pour la référence laminaire} : mesurer $u_\infty$ sur un plateau
    précoce, avant que $k$ n'ait crû --- inutile ici, le biais étant chiffré ;
    \item \textbf{pour un besoin de rigueur accrue} : une boîte plus grande
    ($L_0/D\gtrsim12$ contre $7{,}5$ ici) laisse le sillage se dissiper avant de se
    reboucler --- coûteux ($\times8$ en volume), à réserver à une vérification ponctuelle ;
    \item \textbf{pour les runs turbulents} (l'objet du stage) : le problème
    \emph{se résout de lui-même}, la turbulence forcée décorrélant et brisant le
    sillage à chaque instant --- le sillage propre n'y est qu'une petite perturbation
    sur un champ maintenu stationnaire (cf. la caractérisation \og{}couplage à sens
    unique\fg{} déjà établie).
\end{enumerate}

\begin{figure}[htbp]
    \centering
    \includegraphics[width=0.48\textwidth]{lam7_still_vitesse.png}\hfill
    \includegraphics[width=0.48\textwidth]{lam7_still_vorticite.png}
    \caption{\textbf{Bulle laminaire corrigée} (état terminal, coupe $x$--$y$ à la
    hauteur de la bulle, repère mobile). À gauche : champ de vitesse verticale $u_z$
    (le liquide défile devant la bulle maintenue au centre) ; à droite : module de
    vorticité $|\omega|$, révélant le sillage responsable de l'instabilité de
    trajectoire oblique. Film complet : \texttt{lam7\_vitesse.mp4},
    \texttt{lam7\_vorticite.mp4} (référence).}
    \label{fig:lam7_film}
\end{figure}

\subsection{Convergence en maillage}
\label{subsec:convergence}

La convergence annoncée avant correction --- \og{}lvl7 $\approx$ lvl8 sur $U_b$\fg{} ---
était sans valeur : les deux niveaux s'épinglaient sur la même frontière. Elle a été
refaite proprement.

\paragraph{Pourquoi sur le cas laminaire.} En turbulence, deux maillages diffèrent dès
le premier pas de temps et le chaos amplifie l'écart exponentiellement : les
\emph{trajectoires} divergent nécessairement, même lorsque les \emph{statistiques}
convergent. Une étude de convergence point à point y est donc impossible par
construction. Le cas laminaire, lui, est déterministe : la comparaison est licite, et
elle isole l'effet du maillage de celui du chaos.

\begin{table}[htbp]
    \centering
    \renewcommand{\arraystretch}{1.3}
    \begin{tabular}{lcccccc}
        \toprule
        Niveau & $D/\Delta x$ & $\delta/\Delta x$ & $u_\infty$ & $Re_b$ & $C_D$ & $\chi$ \\
        \midrule
        6 & $8{,}5$  & $0{,}87$ & $10{,}93$  & $95{,}9$  & $0{,}713$ & $1{,}26$ \\
        7 & $17{,}1$ & $1{,}64$ & $12{,}274$ & $107{,}7$ & $0{,}566$ & $1{,}38$ \\
        8 & $34{,}1$ & $3{,}28$ & $12{,}314$ & $108{,}0$ & $0{,}562$ & $1{,}36$ \\
        \bottomrule
    \end{tabular}
    \caption{Convergence en maillage du cas laminaire ($Bo=1$, $Ga=70$, mêmes
    paramètres physiques aux trois niveaux). $\delta\sim D/\sqrt{Re_b}$ est
    l'épaisseur de couche limite. Écart lvl6$\to$lvl7 : $+12{,}3\%$ ;
    \textbf{lvl7$\to$lvl8 : $+0{,}32\%$}. Extrapolation de Richardson ($p=2$) :
    $u_\infty\to12{,}33$, dont lvl7 est à $-0{,}4\%$.}
    \label{tab:convergence}
\end{table}

\paragraph{Ce n'est pas l'interface qui pilote, c'est la couche limite.} Le tableau
le montre : à \texttt{lvl6} la couche limite fait $0{,}87$ maille --- elle n'est pas
résolue --- et $u_\infty$ est $12\%$ trop bas, $C_D$ $27\%$ trop haut. Dès
$\delta/\Delta x\simeq1{,}6$ (\texttt{lvl7}), doubler encore la résolution ne change
plus $u_\infty$ que de $0{,}3\%$. Le comptage usuel \og{}$D/\Delta x$\fg{} est donc
trompeur : ce qui compte est la résolution de la couche limite, et $17$ mailles par
diamètre y suffisent à ce $Re_b$.

\begin{figure}[htbp]
    \centering
    \includegraphics[width=\textwidth]{convergence_maillage.png}
    \caption{\textbf{Convergence en maillage, cas laminaire.} En haut à gauche :
    trajectoires $z_{\text{lab}}(t)$ superposées --- lvl7 et lvl8 sont confondues,
    lvl6 décroche. En haut à droite : $u_\infty$ en fonction de $\delta/\Delta x$,
    le paramètre qui gouverne réellement. En bas à gauche : vitesse instantanée ---
    même transitoire (accélération initiale de masse ajoutée), même cycle limite,
    seul le plateau de lvl6 est décalé.}
    \label{fig:convergence_maillage}
\end{figure}

\begin{table}[htbp]
    \centering
    \renewcommand{\arraystretch}{1.25}
    % genere par scripts/table_resolution.py
\begin{tabular}{lccccccc}
\hline
Point & $\beta$ & $Re_\lambda$ & $\varepsilon$ & $\eta$ & $k_{max}\eta$ & $D/\Delta x$ & $\delta/\Delta x$ \\
\multicolumn{5}{l}{} & \multicolumn{3}{c}{(niveau 7)} \\
\hline
wt05 & 0.15 & 17.8 & 0.35 & 2.03 & 6.82 & 17.1 & 1.69 \\
wt11 & 0.22 & 22.6 & 0.80 & 1.66 & 5.57 & 17.1 & 1.87 \\
wt21 & 0.31 & 26.3 & 1.89 & 1.34 & 4.49 & 17.1 & 2.10 \\
wt25 & 0.33 & 28.4 & 2.25 & 1.28 & 4.30 & 17.1 & 1.89 \\
wt32 & 0.38 & 29.7 & 3.27 & 1.17 & 3.92 & 17.1 & 2.00 \\
\hline
\end{tabular}

    \caption{Résolution des échelles turbulentes au niveau 7, point par point.
    $k_{\max}\eta\ge1{,}5$ est le critère de DNS résolue ; le minimum de la campagne
    vaut $3{,}92$.}
    \label{tab:resol_turb}
\end{table}

\subsection{Convergence en maillage en turbulence : trajectoires divergentes,
statistique convergée}
\label{subsec:conv_turb}

La convergence précédente était sur le cas laminaire, où il n'y a pas de chaos. Un
tuteur avait soulevé, à juste titre, une objection plus profonde : sur les premiers
runs turbulents, les niveaux 7 et 8 donnaient une vitesse d'ascension \emph{instantanée}
quasi identique, oscillation par oscillation. C'est \emph{a priori} surprenant --- un
écoulement turbulent est chaotique, deux maillages différents devraient envoyer la bulle
sur des chemins distincts. Deux réponses, la première étant que cet accord-là était en
partie un artefact.

\paragraph{L'accord parfait d'origine était un épinglage.} Sur ces runs (antérieurs au
correctif, avec \texttt{reduced.h}), les deux bulles étaient bloquées contre la même
frontière $z_c\simeq L_0-R_0$. La corrélation entre leur accélération et leur écart à
cette position vaut $-0{,}42$ : c'est la signature d'un \textbf{rappel élastique}, deux
bulles oscillant dans le même puits de potentiel numérique --- situé en un point du
domaine, indépendant du maillage. Elles battaient donc en phase. Et l'accord n'était que
visuel : l'écart quadratique relatif entre les deux courbes passe de $1{,}4\%$ pendant la
montée libre (dynamique déterministe de flottabilité, accord réel) à $88\%$ pendant la
phase épinglée. La figure~\ref{fig:meme_vitesse} le décompose.

\begin{figure}[htbp]
    \centering
    \includegraphics[width=\textwidth]{pourquoi_meme_vitesse.png}
    \caption{\textbf{Décomposition de l'accord d'origine} (runs $We_t\simeq1$
    antérieurs au correctif). L'accord réel se limite à la montée libre ($t<85$,
    déterministe). Après, les deux bulles sont épinglées au même mur (bas gauche) et y
    oscillent en rappel élastique (bas milieu, pente négative) : leur ressemblance est
    un artefact, l'écart quadratique atteignant $88\%$ (bas droite).}
    \label{fig:meme_vitesse}
\end{figure}

\paragraph{Le vrai test, non épinglé.} On a donc refait trois runs turbulents
\emph{strictement identiques sauf le maillage} (niveaux 6, 7, 8 ; même précurseur
$KE=11$, $\beta=0{,}22$ ; \textbf{même position d'injection} ; pile corrigée). Le
verdict sépare nettement les deux notions que la turbulence mélange
(fig.~\ref{fig:conv_turb} et \ref{fig:traj}) :
\begin{itemize}
  \item les \textbf{trajectoires divergent} : partant du même point, l'écart horizontal
  entre deux maillages croît de ${\sim}0{,}5$ à plus d'un diamètre en $40$~u.t. --- le
  tuteur a raison, les chemins individuels ne sont pas reproductibles ;
  \item mais la \textbf{statistique converge} : $Re_\lambda=22{,}3$ aux trois niveaux, et
  la vitesse d'ascension moyenne vaut $7{,}9$ (lvl6), $11{,}5$ (lvl7) et $11{,}7$ (lvl8)
  sur la fenêtre commune --- soit \textbf{lvl7 $\approx$ lvl8 à $2\%$ près} (dans la
  dispersion de réalisation unique), tandis que \textbf{lvl6 décroche de $30\%$},
  exactement comme en laminaire, sa couche limite ($0{,}87$ maille) et sa turbulence
  ($k_{\max}\eta=2{,}9$) étant toutes deux à la limite.
\end{itemize}
\textbf{La convergence en maillage d'un écoulement turbulent porte donc sur les
statistiques, jamais sur les trajectoires} --- ce qui justifie, et impose, les moyennes
d'ensemble de la section suivante. (Le niveau 8 turbulent était encore en cours au moment
de la rédaction : les valeurs le concernant sont provisoires, sur la fenêtre $[180,209]$
commune aux trois.)

\begin{figure}[htbp]
    \centering
    \includegraphics[width=\textwidth]{convergence_turbulente.png}
    \caption{\textbf{Convergence en maillage en turbulence} ($\beta=0{,}22$, même
    injection). Montée (haut gauche) : lvl7 et lvl8 ensemble, lvl6 en dessous. $u_\infty$
    par niveau (haut droite) : lvl7 $\approx$ lvl8, lvl6 décroche. Trajectoires (bas
    gauche) : elles divergent. Statistiques (bas milieu) : $Re_\lambda$ identique.
    Résolution (bas droite) : Kolmogorov résolu dès lvl7.}
    \label{fig:conv_turb}
\end{figure}

\begin{figure}[htbp]
    \centering
    \includegraphics[width=\textwidth]{trajectoires_cote_a_cote.png}
    \caption{\textbf{Trajectoires de la bulle aux trois maillages, côte à côte.} Même
    turbulence, même point d'injection (cercle noir). Rangée du haut : l'errance
    horizontale (vue de dessus) est \emph{différente} à chaque niveau. En bas : les
    chemins superposés partent ensemble puis se séparent (gauche), la montée garde la même
    pente moyenne (milieu), et la distance entre trajectoires dépasse un diamètre
    (droite). Le chaos décorrèle les chemins ; seule la statistique converge.}
    \label{fig:traj}
\end{figure}

\subsection{Perspective : mesurer enfin \texorpdfstring{$u_\infty$}{u\_inf} et le rapport
\texorpdfstring{$\beta$}{beta}}

Le correctif étant validé et $u_\infty = 12{,}27$ désormais \textbf{mesuré}, la
campagne \texttt{weber\_rise} reprend le protocole dans le cadre du domaine
\cite{legendre2025} : à \textbf{gravité fixe} ($g_0=4$, $Ga\simeq70$, $Bo=1$), le
run laminaire fournit $u_\infty$ et les runs turbulents balayent l'intensité sur les
précurseurs déjà disponibles. Le livrable est la courbe
$\overline{u_\infty}/u_\infty$ en fonction de $\beta = u'/u_\infty$ et du Froude
turbulent $Fr' = u'/\sqrt{gD}$, directement comparable aux figures~24 et~27(d) de la
revue, où la littérature rapporte une \textbf{réduction} de la vitesse de remontée
(${\sim}35\%$ chez Poorte \& Biesheuvel) suivant
$\overline{u_\infty}/u_\infty - 1 \sim \beta^2$ à petit $\beta$, puis
$\overline{u_\infty}/u_\infty \sim 1/Fr'$ à grand $\beta$.

\paragraph{Plage de $\beta$ réellement accessible.} Un point important, que seule la
mesure de $u_\infty$ permet de trancher : avec $u_\infty=12{,}27$ et les cinq
précurseurs disponibles ($u' = \sqrt{2k/3}$ de $1{,}8$ à $4{,}6$), le rapport
$\beta$ ne couvre que $[0{,}15;\,0{,}38]$ (et $Fr'\in[0{,}23;\,0{,}58]$), et non la
fourchette $[0{,}4;\,0{,}9]$ initialement anticipée --- qui supposait implicitement
un $u_\infty$ bien plus faible. Ces cinq points explorent donc le régime de
\textbf{petit $\beta$}, celui de la loi
$\overline{u_\infty}/u_\infty - 1 \sim \beta^2$.

\paragraph{Cette limitation n'est pas intrinsèque.} Il serait tentant de la justifier
en remarquant qu'atteindre $\beta\simeq0{,}9$ exige $u'\simeq11$, soit $k\simeq180$,
c'est-à-dire le champ $KE\simeq200$ qui \textbf{fragmentait} la bulle --- et d'en
conclure à une tension intrinsèque entre grands $\beta$ et intégrité de la bulle.
\textbf{L'argument ne tient pas}, et il est instructif de voir pourquoi. Ce run de
fragmentation était à $g=1$, donc à $\sigma=(\rho_1-\rho_2)gD^2/Bo=255{,}7$. La
campagne corrigée tourne à $g=4$, donc à $\sigma=1022{,}8$ : \textbf{quatre fois
plus}. Or $We_t=\rho u'^2D/\sigma$ est \emph{inversement} proportionnel à $\sigma$ :
à turbulence égale, le Weber turbulent est aujourd'hui quatre fois plus petit
qu'alors. En prenant $We_c=3$ \cite{riviere2021}, le seuil de rupture n'est atteint
que vers $u'\simeq13{,}9$, soit $\beta\simeq1{,}1$. La contrainte de résolution
($k_{\max}\eta\ge1{,}5$) ne mord, elle, qu'au-delà. \textbf{C'est donc la
fragmentation qui borne, et elle borne vers $\beta\simeq1$, non vers $0{,}4$} : la
plage accessible peut être triplée sans changer ni $Bo$ ni $Ga$. C'est l'objet de la
campagne décrite au \S\ref{subsec:campagne}.

\subsection{Premiers résultats : la turbulence ralentit la remontée}
\label{subsec:resultats_turb}

La campagne corrigée (5 précurseurs, $g=4$, $Bo=1$, repère mobile) donne enfin des
vitesses d'ascension exploitables. La vitesse terminale turbulente
$\overline{u}_{\infty,\text{turb}}$ est le plateau de \texttt{frame\_uz} moyenné
sur $t\ge220$ ; on la rapporte à la référence laminaire $u_\infty=12{,}27$
(tab.~\ref{tab:resultats_turb}, fig.~\ref{fig:weber_rise}).

\begin{table}[htbp]
    \centering
    \renewcommand{\arraystretch}{1.3}
    \begin{tabular}{lccccc}
        \toprule
        Run & précurseur & $\beta=u'/u_\infty$ & $\overline{u}_{\infty,\text{turb}}$ &
        $\overline{u}_{\infty,\text{turb}}/u_\infty$ & ralentissement \\
        \midrule
        (laminaire) & --- & $0$ & $12{,}27$ & $1{,}00$ & référence \\
        wt05 & we03 & $0{,}15$ & $11{,}57$ & $0{,}94$ & $-6\%$ \\
        wt11 & we07 & $0{,}22$ & $10{,}09$ & $0{,}82$ & $-18\%$ \\
        wt21 & we13 & $0{,}31$ & $7{,}68$ & $0{,}63$ & $-37\%$ \\
        wt25 & we16 & $0{,}33$ & $8{,}35$ & $0{,}68$ & $-32\%$ \\
        wt32 & we20 & $0{,}38$ & $7{,}53$ & $0{,}61$ & $-39\%$ \\
        \bottomrule
    \end{tabular}
    \caption{Vitesse d'ascension en fonction de l'intensité turbulente
    (\textbf{moyenne d'ensemble} par point, injections à positions décalées dans le
    même précurseur : $n=5$ réalisations pour $\beta=0{,}15$, $n=3$ pour les quatre
    autres, dont les membres manquants sont en cours de calcul).
    \textbf{La turbulence ralentit systématiquement la remontée}, d'abord vite, puis
    beaucoup plus lentement au-delà de $\beta\simeq0{,}3$.}
    \label{tab:resultats_turb}
\end{table}

\begin{figure}[htbp]
    \centering
    \includegraphics[width=\textwidth]{courbe_maitresse.png}
    \caption{\textbf{Ralentissement de la remontée par la turbulence} ($Bo=1$,
    $Ga=70$, niveau 7 validé en maillage). Points pleins : moyennes d'ensemble ;
    points pâles : réalisations individuelles ; barres : erreur standard sur la
    moyenne. \textbf{À gauche}, en fonction de $\beta=u'/u_\infty$, avec la loi
    petit-$\beta$ ajustée. \textbf{À droite}, en fonction du Froude turbulent
    $Fr'=u'/\sqrt{gD}$, superposée aux deux branches de la littérature : nos points
    quittent la parabole vers $Fr'\simeq0{,}45$ et s'infléchissent vers la loi
    $0{,}37/Fr'$ --- ils échantillonnent la \emph{transition} entre les deux régimes.}
    \label{fig:weber_rise}
\end{figure}

\paragraph{Trois enseignements.} \textbf{(1)} La turbulence \emph{ralentit}
systématiquement la bulle, d'abord rapidement ($-6\%$ à $\beta=0{,}15$, $-18\%$ à
$\beta=0{,}22$, $-37\%$ à $\beta=0{,}31$), puis \textbf{beaucoup plus lentement} :
les trois points $\beta\ge0{,}3$ ($-37\%$, $-32\%$ et $-39\%$) sont
\emph{indiscernables entre eux} compte tenu des barres d'erreur (SEM de $1{,}2$ à
$4{,}5$ points). L'échantillon actuel ne permet pas de trancher entre un véritable
plateau et une décroissance lente de type $1/Fr'$ ; il exclut en revanche la
poursuite de la loi quadratique, qui prédirait $-53\%$ à $\beta=0{,}38$, soit trois
écarts-types sous la mesure. La bulle n'est \textbf{jamais fragmentée} : le contrôle
recalculé sur les $17$ réalisations avec le diagnostic corrigé donne un volume
$\ge97{,}1\%$ du théorique et une \emph{seule} région gazeuse à tout instant, pour un
rapport d'aspect $\chi\le1{,}67$.
\textbf{(2)} La réduction dépasse nettement la loi petit-$\beta$ dès
$\beta\simeq0{,}2$, ce qui place l'écoulement dans le régime de forte hindrance
($\sim1/Fr'$) plutôt que quadratique --- l'ordre de grandeur rejoint les $-35\%$ de
Poorte \& Biesheuvel.
\textbf{(3) Nécessité de la moyenne d'ensemble} : à $\beta\gtrsim0{,}3$,
$u'/\overline{u}_\infty\simeq0{,}5$ et une \emph{seule} réalisation ne fixe pas la
moyenne. Une première mesure isolée donnait d'ailleurs une non-monotonie trompeuse
($-24\%$ à $\beta=0{,}33$ contre $-37\%$ à $\beta=0{,}31$) ; en moyennant trois
injections indépendantes, ce point remonte à $-32\%$ et l'anomalie disparaît.

\subsection{Campagne en cours : étendre et fiabiliser la courbe}
\label{subsec:campagne}

Un premier balayage dense (neuf points, $\beta$ de $0{,}05$ à $0{,}45$) avait été
soumis, puis \emph{annulé} avant d'avoir produit le moindre résultat exploitable. Il
a été remplacé le 24 juillet par une campagne de \textbf{cinquante travaux} qui
attaque les trois faiblesses de la figure~\ref{fig:weber_rise}.

\begin{description}
  \item[(A) $n=5$ partout.] Deux réalisations supplémentaires pour $\beta=0{,}22$ ;
  $0{,}31$ ; $0{,}33$ ; $0{,}38$ (le point $\beta=0{,}15$ les a déjà). L'argument
  n'est pas la taille des barres d'erreur --- passer de $3$ à $5$ ne divise la SEM
  que par $1{,}3$ --- mais leur \emph{fiabilité} : à $n=3$, l'incertitude
  \emph{sur l'écart-type lui-même} vaut $1/\sqrt{2(n-1)}\simeq50\%$. Annoncer
  \og{}$\pm4{,}5\%$\fg{} quand la valeur réelle peut être $2\%$ ou $7\%$ n'a pas de
  sens ; à $n=5$ cette incertitude retombe à ${\sim}35\%$. Le point $\beta=0{,}15$
  illustre le risque : son cinquième membre est sorti à $10{,}24$ quand les quatre
  autres se tenaient entre $11{,}6$ et $12{,}2$ --- une dispersion que $n=3$ aurait
  masquée.
  \item[(B) Trois points bas] ($\beta = 0{,}050$ ; $0{,}075$ ; $0{,}100$). La loi
  $1-c\,\beta^2$ est aujourd'hui ajustée sur des points tous $\ge0{,}15$, ancrée par
  le seul point laminaire : le comportement quadratique à l'origine est
  \textbf{supposé, jamais vérifié}. Ces trois points le testent directement.
  \item[(C) Quatre points hauts] ($\beta = 0{,}50$ ; $0{,}65$ ; $0{,}85$ ; $1{,}00$),
  qui \textbf{triplent la plage} explorée. Comme montré plus haut, la fragmentation
  ne borne qu'à $\beta\simeq1{,}1$ : $We_t$ n'atteint que $2{,}35$ au point le plus
  intense, sous le seuil $We_c=3$. C'est dans cette plage que la branche $0{,}37/Fr'$
  domine et que la confrontation à la littérature devient discriminante. Et si la
  bulle fragmentait malgré tout, ce ne serait pas un échec mais une \textbf{mesure de
  $We_c$} : l'expérience est concluante dans les deux cas.
\end{description}

\begin{figure}[htbp]
    \centering
    \includegraphics[width=\textwidth]{campagne_beta.png}
    \caption{\textbf{Conception de la campagne.} À gauche : couverture en $\beta$,
    des cinq points acquis (pleins) aux douze visés (cercles ouverts, placés à leur
    position \emph{attendue} sur la branche pertinente). À droite : $We_t$ point par
    point --- c'est parce que la campagne tourne à $g=4$, donc à $\sigma=1023$, que
    l'on peut monter jusqu'à $\beta=1$ sans franchir le seuil de fragmentation.}
    \label{fig:campagne}
\end{figure}

\subsection{Confrontation à la littérature : le cadre de Froude}
\label{subsec:froude}

Notre paramètre de balayage $\beta = u'/u_\infty$ est équivalent au \textbf{Froude
turbulent} $Fr' = u'/\sqrt{gD}$ employé dans la littérature récente, les deux étant
reliés par le Froude d'ascension quiescent $u_\infty/\sqrt{gD}=1{,}53$
(soit $Fr'=1{,}53\,\beta$). Or l'étude DNS de Liu \emph{et al.}~\cite{liu2024}
--- qui traite \emph{exactement} notre problème (bulle montant en THI, même solveur
Basilisk, même protocole précurseur, $\rho_2/\rho_1=1/850$) --- établit que la
réduction de vitesse d'ascension est \textbf{principalement fonction de $Fr'$}, les
nombres $Ga$ et $Bo$ n'ayant qu'un rôle secondaire sur la réduction \emph{relative}.
Deux régimes se dégagent :
\begin{equation}
    \frac{\overline{u}_{\infty,\text{turb}}}{u_\infty} = 1 - \chi\,Fr'^2
    \quad (Fr'\lesssim0{,}5), \qquad
    \frac{\overline{u}_{\infty,\text{turb}}}{u_\infty} = \frac{0{,}37}{Fr'}
    \quad (Fr'\gtrsim0{,}5),
    \label{eq:froude}
\end{equation}
le régime de grand $Fr'$ suivant la loi $\propto 1/Fr'$ de Ruth et al. (2021).

La figure~\ref{fig:froude} superpose nos cinq points (moyenne d'ensemble, lvl7) à
ces deux lois. \textbf{Nos résultats tombent sur la courbe maître} : le point de plus
faible intensité ($Fr'=0{,}23$) suit la parabole de Spelt \& Biesheuvel, et les points
de forte intensité ($Fr'\gtrsim0{,}5$) rejoignent la loi $0{,}37/Fr'$ (p.~ex.
$Fr'=0{,}58$ : mesuré $0{,}61\pm0{,}04$ contre $0{,}64$, soit ${\sim}4\%$ d'écart).
L'accord, obtenu \emph{sans
ajustement}, valide à la fois la chaîne de calcul corrigée et le fait que $Fr'$ est
le bon paramètre de contrôle.

\begin{figure}[htbp]
    \centering
    \includegraphics[width=0.72\textwidth]{weber_vs_Fr_litterature.png}
    \caption{Vitesse d'ascension relative $\overline{u}_{\infty,\text{turb}}/u_\infty$
    en fonction du Froude turbulent $Fr'=u'/\sqrt{gD}$. Points rouges : nos DNS
    (moyenne d'ensemble, lvl7) ; carré noir : référence laminaire ($Fr'=0$). Tirets
    bleus : loi petit-$Fr'$ $1-\chi Fr'^2$ \cite{spelt1997} ; trait vert :
    loi grand-$Fr'$ $0{,}37/Fr'$ \cite{ruth2021,liu2024}. Nos points
    tracent la transition entre les deux régimes.}
    \label{fig:froude}
\end{figure}

\subsection{Identification du mécanisme : traînée non linéaire, mesurée}
\label{subsec:mecanisme}

Deux mécanismes peuvent expliquer le ralentissement, et ils ont des signatures
\emph{opposées}. \textbf{(A) Échantillonnage préférentiel} : la bulle passerait plus de
temps dans du fluide qui descend, et serait ainsi \emph{portée} vers le bas ; sa
signature est une vitesse verticale du fluide vu par la bulle nettement \emph{négative},
du même ordre que la réduction $\Delta U$. \textbf{(B) Traînée non linéaire}
\cite{liu2024} : le fluide ambiant a une vitesse moyenne nulle, mais $C_D(Re)$ étant
non linéaire, les fluctuations de la vitesse de glissement augmentent la traînée
\emph{moyenne} ; sa signature est un fluide ambiant $\simeq0$ mais une forte
\emph{variance} du glissement.

\paragraph{Mesure.} On échantillonne la vitesse verticale du fluide dans un cône
supérieur autour de la bulle (\texttt{src/mechanism.c}, snapshots des runs
$\beta=0{,}15$, $0{,}31$, $0{,}38$), à plusieurs rayons, en retranchant l'écoulement que
la bulle \emph{induit} elle-même --- calibré sur le run laminaire, où il n'existe aucune
turbulence à échantillonner (il y vaut $+0{,}24$ à $r=3D$, décroissant, et positif comme
attendu en amont d'une bulle montante). La différence proche$-$lointain est indépendante
du repère mobile (le terme $\mathrm{frame\_uz}$ s'y annule). Le tableau ci-dessous oppose
la réduction totale $\Delta U = u_\infty - \overline{W}$ à la part qu'expliquerait
l'échantillonnage préférentiel.

\begin{center}
\renewcommand{\arraystretch}{1.25}
\begin{tabular}{lcccc}
    \toprule
    Run & $\beta$ & $\Delta U$ & signal préférentiel & part de $\Delta U$ \\
    \midrule
    wt05 & $0{,}15$ & $0{,}66$ & $-0{,}01$ & $1\%$ \\
    wt21 & $0{,}31$ & $4{,}71$ & $-0{,}57$ & $12\%$ \\
    wt32 & $0{,}38$ & $3{,}96$ & $-0{,}09$ & $2\%$ \\
    \bottomrule
\end{tabular}
\end{center}

\paragraph{Verdict.} Une fois l'écoulement induit retranché, le fluide vu par la bulle
est \textbf{quasi nul} : l'échantillonnage préférentiel explique au plus $12\%$ de la
réduction, et ce point isolé est vraisemblablement du bruit (la tendance n'est pas
monotone --- $\beta=0{,}38$, pourtant plus intense, ne donne que $2\%$). En regard, les
\textbf{fluctuations de la vitesse de glissement} croissent de $8\%$ à $50\%$ de
$\beta=0{,}15$ à $0{,}38$ (fig.~\ref{fig:mecanisme}) : c'est exactement l'amplitude que
requiert la traînée non linéaire. \textbf{Le ralentissement est donc dominé par la
traînée non linéaire}, en accord avec \cite{liu2024} pour une bulle super-Kolmogorov
($St\gg1$), l'échantillonnage préférentiel --- réservé aux bulles sous-Kolmogorov
($St\sim1$) --- restant marginal. On assume les limites de cette mesure : réalisations
uniques, coquille située à l'intérieur de l'échelle intégrale (approximation), trois
points ; mais le signe et l'ordre de grandeur sont robustes.

\begin{figure}[htbp]
    \centering
    \includegraphics[width=\textwidth]{mechanism.png}
    \caption{\textbf{Mécanisme du ralentissement.} À gauche : vitesse du fluide vue par
    la bulle selon le rayon ; le laminaire (gris) donne l'écoulement induit à
    retrancher. Au centre : la réduction totale (rouge) contre la part expliquée par
    l'échantillonnage préférentiel (violet) --- l'écart est la traînée non linéaire. À
    droite : les fluctuations de glissement croissent avec $\beta$, condition de la
    traînée non linéaire.}
    \label{fig:mecanisme}
\end{figure}

\paragraph{Cohérence avec la trajectoire.} Ce diagnostic est cohérent avec ce que l'on
observe :
la bulle suit une trajectoire \textbf{oblique quasi-rectiligne} (déviation ${\sim}5$--$7^\circ$
de la verticale, $Ga=70$ étant juste au-dessus du seuil $Ga_c\simeq44$), le long de
laquelle la vitesse d'ascension présente une \textbf{modulation périodique faible} :
sur le plateau, $w$ fluctue avec un écart-type de $0{,}21$, soit $1{,}7\%$ de
$u_\infty$, à la période du lâcher tourbillonnaire ($T\simeq16{,}2$~u.t.,
$St=0{,}080$). C'est une modulation du glissement bien moins spectaculaire que la
spirale du cas $Ga=237$ de Liu et al., mais relevant du même couplage
bulle--sillage.\footnote{Une version antérieure de ce document décrivait ici une
bulle qui \og{}rampe\fg{}. Cette observation est \textbf{retirée} : elle provenait
d'un artefact de rendu. Les films étaient produits par une coupe \emph{horizontale}
($n=\{0,0,1\}$, caméra \texttt{front}), qui place l'axe d'ascension perpendiculaire à
l'écran et rend la montée géométriquement invisible ; seule subsistait à l'image
l'oscillation latérale. Corrigé dans \texttt{src/render\_lab.c} (coupe verticale
$n=\{0,1,0\}$, caméra \texttt{bottom}). La bulle monte en réalité de façon
parfaitement régulière, à $97{,}6\%$ de $u_\infty$ au minimum.}
Le seuil de non-fragmentation est lui aussi cohérent : nos $We_t\le0{,}34$ sont bien
sous le $We_c=3$ mesuré par Rivière et al. (2021), d'où l'absence de rupture.

\paragraph{Réserve de résolution : levée.} Nos points couvrent
$Fr'\in[0{,}23;0{,}58]$, comblant le bas de la gamme de Liu et al. Une différence de
maillage doit être signalée : leurs DNS sont à ${\sim}68$ mailles par diamètre,
contre ${\sim}17$ pour notre niveau 7. Ce comptage brut est toutefois trompeur, et
l'étude de convergence du \S\ref{subsec:convergence} montre pourquoi : ce qui pilote
la vitesse d'ascension n'est pas le nombre de mailles sur le diamètre, mais la
résolution de la \textbf{couche limite} $\delta\sim D/\sqrt{Re_b}$. Le saut
lvl6$\to$lvl7 ($+12{,}3\%$) correspond exactement au passage d'une couche limite non
résolue ($\delta/\Delta x=0{,}87$) à une couche limite résolue ($1{,}64$) ; au-delà,
doubler encore la résolution ne déplace $u_\infty$ que de $0{,}32\%$. Le niveau 7 est
donc \textbf{convergé pour la grandeur mesurée}, et l'accord avec Liu et al.\ n'est
pas une coïncidence de résolutions différentes.

\paragraph{Leçon de méthode.} C'est le run \emph{laminaire} --- dont la seule
raison d'être était de fournir une référence $u_\infty$ --- qui a révélé
l'artefact. Sans lui, la campagne aurait conclu que \og{}la turbulence supprime la
remontée de la bulle\fg{}, résultat spectaculaire, cohérent avec la littérature
\cite{legendre2025} \emph{et entièrement faux}. Un cas limite dont la réponse est
connue \emph{a priori} (ici : une bulle doit monter indéfiniment en fluide au
repos) est le meilleur détecteur d'artefact numérique.
